\documentclass[10pt,a4paper]{amsart}
\usepackage[margin=19mm]{geometry}
\usepackage{amsmath,amssymb,mathtools}
\usepackage[scr=rsfs,cal=euler]{mathalfa}
\usepackage{xfrac}
\usepackage[unicode,hidelinks,bookmarks=false]{hyperref}
\newcommand{\ie}{i.e.\ }
\newcommand{\cf}{cf.\ }
\newcommand{\resp}{respectively\ }
\newcommand{\Subsection}{\subsection}
\providecommand{\nobreakdash}{\nobreak\hbox{-}}
\setlength{\emergencystretch}{2em}
\multlinegap0pt
\usepackage{tikz-cd}
\usepackage{enumitem}
\tikzcdset{arrow style=tikz}
\numberwithin{equation}{section}
\renewcommand{\mathbb}{\mathbf}
\renewcommand{\setminus}{\mathbin{\rule[0.2em]{0.67em}{0.12em}}}
\allowdisplaybreaks
\theoremstyle{plain}
\newtheorem{theorem}[equation]{Theorem}
\newtheorem{lemma}[equation]{Lemma}
\theoremstyle{definition}
\newtheorem{situation}[equation]{}
\theoremstyle{plain}
\setcounter{section}{2}
\setcounter{equation}{1}
\title{Revisiting Dwork cohomology: visibility and divisibility of Frobenius eigenvalues in rigid cohomology: exactness of the rows of the Dwork double complex (Lemma 6.7)}
\author{Daqing Wan}
\author{Dingxin Zhang}
\date{Source: arXiv:2207.12633v2; \url{https://arxiv.org/abs/2207.12633v2}}
\begin{document}
\maketitle
\noindent\textit{Editorial scope.} Counted unit: the exactness of the rows of the article's Dwork double complex, printed as Lemma 6.7 (source0.tex lines 2570--2574, source label \texttt{lemma:c-exactness-of-rows}), delivered together with the article's own complete proof of it (source0.tex lines 2576--2710, the adjacent proof block that begins with the change from the toric to the affine convention and ends with ``This completes the proof.''). No line of the counted statement or of the counted proof is added, dropped, reordered or re-flowed, no mathematics is written, repaired or re-derived, and the article's printed slips are kept literally, for instance the missing prime in the module $\mathscr{D}^{\dagger}_{\mathcal{P}^{\prime},\mathbb{Q}}({}^{\dagger}H)$ inside the proof, where the surrounding text writes ${}^{\dagger}H^{\prime}$.
Required context retained uncounted, each block being either cited by the counted proof or defining an object the counted statement names: the article's Situation 2.2 that the comparison theorem refers to (lines 604--613); the Baldassarri--Berthelot comparison theorem as the article states it, printed as Theorem 2.18 (lines 920--930), which the proof invokes twice; the article's lemma on the concentration of the local cohomology complex, printed as Lemma 5.6 (lines 2318--2328), which the proof invokes once together with the retained hypothesis; the article's additional hypothesis used in the proof, printed as equation (6.4) (lines 2529--2531); and the article's own definition of the double complex, printed as equation (6.5), together with the differentials (6.6) and the sentence reading them off (lines 2548--2568), which the counted statement quantifies over. Every retained context block is the article's text line for line. The sentence that first introduces the double complex in the source is not retained because it refers back to the article's earlier chain-level construction; the complex itself, its differentials and the sentence that identifies them are retained, so the counted statement is complete as to its own objects.
Editorial formatting only, in addition: an AMS \texttt{amsart} wrapper that restates the article's own environment declarations (the \texttt{theorem}, \texttt{lemma} and \texttt{situation} environments sharing the single equation counter exactly as the article declares them, with equations numbered within the section) and the article's own typographic conventions that the retained lines rely on (the redefinition of \texttt{\textbackslash mathbb} as bold and of \texttt{\textbackslash setminus} as a short rule, and the \texttt{tikz-cd} package for the two diagrams, drawn with the article's arrow style); and six counter settings that reproduce the article's own numbering of the retained blocks, namely Situation 2.2, Theorem 2.18, Lemma 5.6, equations (6.4), (6.5), (6.6), Lemma 6.7 and equation (6.8), with the tagged display (6.5') carrying its own tag as in the article. The article is typeset in a publisher-specific format with its own page geometry and fonts, so the visual presentation differs although the mathematical text does not.
Bibliography: the retained text cites exactly two entries of the article's own bibliography, and both are transcribed verbatim from the article's \texttt{.bbl} file with the \texttt{\textbackslash newblock} line breaks of that file removed and with the article's own printed numbers [8] and [11] as labels; no other entry of the article's bibliography is reproduced and no citation outside those two occurs in any retained line. No figure is retained: no retained line contains a graphic inclusion, and the wrapper declares no asset.
\begin{situation}\label{situation}%
Let \(f_{1},\ldots,f_{r} \in \mathbb{F}_{q}[x_{1},\ldots,x_{n}]\).
Let
\(Z=\operatorname{Spec}(\mathbb{F}_{q}[x_{1},\ldots,x_{n}]/(f_{1},\ldots,f_{r}))\).
Introducing new variables \(x_{n+1},\ldots,x_{n+r}\), let
\(g = x_{n+1}f_{1} + \cdots + x_{n+r}f_{r}\).  We shall refer to the
rigid cohomology with twisted coefficient
\(\mathrm{H}^{\bullet}_{\mathrm{rig}}(\mathbb{A}_{\mathbb{F}_{q}}^{n+r};g^{\ast}\mathcal{L}_{\pi})\),
as the ``overconvergent Dwork cohomology'' of \(Z\).
\end{situation}

\setcounter{equation}{17}
\begin{theorem}[{\cite[Theorem~2.14]{baldassarri-berthelot:dwork-cohomology-for-singular-hypersurfaces}}]%
\label{theorem:local-comparison}
Let notation be as in Situation~\ref{situation}.  Let
\(\varpi\colon \mathbb{A}^{n+r} \to \mathbb{A}^{n}\) be the projection
to the first \(n\) coordinates.  Let \(\mathcal{L}\) denote the
arithmetic \(\mathscr{D}\)-module associated to the Dwork crystal
\(g^{\ast}\mathcal{L}_{\pi}\) on \(\mathbb{A}^{n+r}\).  Then \(\varpi_{+}\mathcal{L}\) is
isomorphic to the local cohomology complex
\(\mathbb{R}\underline{\Gamma}_{Z}^{\dagger}(\mathcal{O}_{\widehat{\mathbb{P}}^{n},\mathbb{Q}}({}^{\dagger}H))[r]\).
% \textup{(See Theorem~\ref{theorem:local-comparison-1}.)}
\end{theorem}

\setcounter{section}{5}\setcounter{equation}{5}
\begin{lemma}\label{lemma:local-cohomology}
Regard \(\mathbb{A}^{N}_{k}\) as a locally closed subscheme of the
formal projective space \(\widehat{\mathbb{P}}^{N}\) over
\(\operatorname{Spf}(\mathcal{O}_{K})\).  Let \(f_{1},\ldots, f_{r}\)
be regular functions on \(\mathbb{A}^{N}_{k}\), defining a closed
subscheme \(Z\) of \(\mathbb{A}^{N}_{k}\).  Assume that
\(\dim Z = N-r\), then the local cohomology
\(\mathscr{D}^{\dagger}_{\widehat{\mathbb{P}}^{N},\mathbb{Q}}\)-modules
\(\mathcal{H}^{m}\{\mathbb{R}\underline{\Gamma}_{Z}^{\dagger}(\mathcal{O}_{\widehat{\mathbb{P}}^{N},\mathbb{Q}}({}^{\dagger}H))[r]\}\)
are zero unless \(m = 0\).
\end{lemma}

\setcounter{section}{6}\setcounter{equation}{3}
\begin{equation}\label{eq:hypothesis}
  \textit{The scheme \(\operatorname{Spec}\mathbb{F}_q[x_1,\ldots,x_n]/(f_1,\ldots,f_c)\) has dimension equal to \(\dim Z\).}
\end{equation}

\begin{equation}\label{eq:double-complex}
  \begin{tikzcd}[column sep=small]
    \bigoplus\limits_{\substack{|I^{\prime}|=n \\|I^{\prime\prime}|=0}}B_I \ar[r] & \cdots \ar[r] & \bigoplus\limits_{\substack{|I^{\prime}|=n \\|I^{\prime\prime}|=c}}B_I  \ar[r] & \cdots \ar[r] & \bigoplus\limits_{\substack{|I^{\prime}|=n \\|I^{\prime\prime}|=r}}B_I \\
    {\vdots} \ar[u] & {} & {\vdots} \ar[u]  & {} & {\vdots} \ar[u] \\
    \bigoplus\limits_{\substack{|I^{\prime}|=0 \\|I^{\prime\prime}|=0}}B_I \ar[r] \ar[u]& \cdots \ar[r] & \bigoplus\limits_{\substack{|I^{\prime}|=0 \\|I^{\prime\prime}|=c}}B_I \ar[u]  \ar[r] & \cdots \ar[r] & \bigoplus\limits_{\substack{|I^{\prime}|=0 \\|I^{\prime\prime}|=r}}B_I \ar[u]
  \end{tikzcd}
\end{equation}
In the diagram, the horizontal differentials are induced by
\(D^{\prime}_{n+1},\ldots,D^{\prime}_{n+r}\),
and the vertical ones are induced by \(D^{\prime}_1,\ldots,D^{\prime}_{n}\),
where
\begin{equation}
D^{\prime}_{i} = x_{i}\frac{\partial}{\partial x_i} + \pi x_{i}\frac{\partial G}{\partial x_{i}}
=\exp(-\pi G) \circ x_{i}\frac{\partial}{\partial x_i}\circ \exp(\pi G), \quad i=1,2,\ldots,n+r.
\end{equation}


The following lemma shows that the \(0\)\textsuperscript{th},
\(1\)\textsuperscript{st}, \(\ldots\), and \((c-1)\)\textsuperscript{st} columns
of the \(E_1\)-page of the spectral sequence associated to the double
complex~\eqref{eq:double-complex} are all zero.

\begin{lemma}
  \label{lemma:c-exactness-of-rows}
  For each \(0\leq i\leq n\), the \(i\)\textsuperscript{th} row of
  \eqref{eq:double-complex} is exact in cohomology degree \(0,1,\ldots,c-1\).
\end{lemma}

\begin{proof}
The complex \eqref{eq:double-complex} uses ``toric'' conventions, and
is more suitable for later chain level manipulations.  In the
following proof, we shall use an equivalent ``affine'' convention.
Write \(B = K\langle x_{1},\ldots,x_{n+r} \rangle^{\dagger}\), and let
\begin{equation*}
D_{j} = \frac{\partial}{\partial x_{j}} + \pi \frac{\partial G}{\partial x_{j}}, \quad j=1,\ldots,n+r.
\end{equation*}
Then \eqref{eq:double-complex} can be written as
\begin{equation}\tag{\ref*{eq:double-complex}\({}^{\prime}\)}
\label{eq:double-complex-prime}
\begin{tikzcd}
\vdots & \vdots & \vdots \\
B^{\oplus n} \ar{r} \ar{u} & B^{\oplus nr} \ar{u} \ar{r} & B^{\oplus n\binom{r}{2}} \ar{r} \ar{u} & \cdots \\
B \ar{r} \ar{u} & B^{\oplus r} \ar{r} \ar{u} & B^{\binom{r}{2}} \ar{r} \ar{u} & \cdots
\end{tikzcd}
\end{equation}
In this double complex, the horizontal arrows are induced by
\(D_{n+1},\ldots,D_{n+r}\), and the vertical arrows are induced by
\(D_{1},\ldots,D_{n}\).

Let us first show the zeroth row of \eqref{eq:double-complex-prime} is
exact in degrees \(0,1,\ldots,c-1\).  Since the
\(i\)\textsuperscript{th} row of \eqref{eq:double-complex-prime} is an
\(\binom{n}{i}\)-fold direct sum of the zeroth row, the desired
exactness in general will follow.

The zeroth row of \eqref{eq:double-complex-prime} is itself the total
complex of a double complex:
\begin{equation}\label{eq:smaller}
\begin{tikzcd}
\vdots & \vdots & \vdots \\
B^{\oplus \binom{r-c}{2}} \ar{r} \ar{u} & B^{\oplus c\binom{r-c}{2}} \ar{r} \ar{u} & B^{\oplus\binom{r-c}{2}\binom{c}{2}} \ar{r}\ar{u} &\cdots \\
B^{\oplus (r-c)} \ar{r} \ar{u} & B^{\oplus c(r-c)} \ar{r}\ar{u} & B^{\oplus(r-c)\binom{c}{2}} \ar{r} \ar{u} &\cdots & \\
B \ar{r} \ar{u} & B^{\oplus c} \ar{r} \ar{u} & B^{\oplus \binom{c}{2}} \ar{r} \ar{u} & \cdots
\end{tikzcd}.
\end{equation}
In the diagram, the horizontal arrows are induced by
\begin{equation*}
D_{n+1},\ldots, D_{n+c},
\end{equation*}
and the vertical arrows are induced by
\begin{equation*}
D_{n+c+1},\ldots, D_{n+r}.
\end{equation*}
Thus the \(i\)\textsuperscript{th} row is the \(\binom{r-c}{i}\)-fold
direct sum of the zeroth row.  If we proved the zeroth row of
\eqref{eq:smaller} is acyclic except in top cohomology degree, then
every row of \eqref{eq:smaller} will be exact except in top cohomology
degree.  Thus the total complex of \eqref{eq:smaller}, which is the
zeroth row of the original \eqref{eq:double-complex-prime}, will have
vanishing cohomology in degrees \(0,1,\ldots,c-1\).


\medskip%
To prove acyclicity, we make the following auxiliary construction.
Consider the projection
\[
\varpi^{\prime}\colon \mathbb{A}^{n+r} \to \mathbb{A}^{n+r-c}, \quad (x_{1}, \ldots, x_{n+r}) \mapsto (x_{1}, \ldots, x_{n}, x_{n+c+1}, \ldots, x_{n+r}).
\]
Define
\(\mathcal{P} = \widehat{\mathbb{P}}^{n+r-c}_{\mathcal{O}_{K}} \times \widehat{\mathbb{P}}^{c}_{\mathcal{O}_{K}}\)
and
\(\mathcal{P}^{\prime} = \widehat{\mathbb{P}}^{n+r-c}_{\mathcal{O}_{K}}\).
The special fiber of \(\mathcal{P}\) (resp. \(\mathcal{P}^{\prime}\))
contains \(\mathbb{A}^{n+r}\) (resp. \(\mathbb{A}^{n+r-c}\)) as a
Zariski open subset, with complement \(H\) (resp. \(H^{\prime}\)).
Let \(\mathcal{L}\) be the holonomic
\(\mathscr{D}^{\dagger}_{\mathcal{P}, \mathbb{Q}}\)-module associated
with the Dwork crystal \(g^{\prime\ast}\mathcal{L}_{\pi}\), where
\[
g^{\prime} = x_{n+1} f_{1} + \cdots + x_{n+c} f_{c}.
\]
By construction, \(\mathcal{L}\) has an overconvergent singularity
along \(H\), i.e., \(\mathcal{L} = \mathcal{L}({}^{\dagger}H)\), and
is naturally a
\(\mathscr{D}^{\dagger}_{\mathcal{P},\mathbb{Q}}({}^{\dagger}H)\)-module,
where
\(\mathscr{D}^{\dagger}_{\mathcal{P},\mathbb{Q}}({}^{\dagger}H)\) is
the sheaf of differential operators on \(\mathcal{P}\) of finite level
with overconvergent singularities along \(H\)
(see \cite[\S4.2.5]{berthelot:arithmetic-d-modules-1-differential-operators-of-finite-level} or
\cite[\S2.5]{baldassarri-berthelot:dwork-cohomology-for-singular-hypersurfaces}).

Thus, \(\mathcal{L}\) represents an object of
\(D^{b}_{\mathrm{hol}}(\mathbb{A}^{n+r}/K)\). By
Theorem~\ref{theorem:local-comparison}, we have
\begin{equation*}
\varpi^{\prime}_{+} \mathcal{L} \simeq \mathbb{R}\underline{\Gamma}_{Z^{\prime}}(\mathcal{O}_{\mathcal{P}^{\prime}, \mathbb{Q}}({}^{\dagger}H^{\prime}))[c],
\end{equation*}
where \(Z^{\prime}\) is the zero locus of \(f_{1}, \ldots, f_{c}\) in
the affine space \(\mathbb{A}^{n} \times \mathbb{A}^{r-c}\),
that is
\[
Z^{\prime} = \{x \in \mathbb{A}^{n} : f_{1}(x) = \cdots = f_{c}(x) = 0\} \times \mathbb{A}^{r-c}.
\]
Given Hypothesis \eqref{eq:hypothesis}, we have
\(\dim Z^{\prime} = n + r - 2c\), and it follows from
Lemma~\ref{lemma:local-cohomology} that
\(\mathbb{R}\underline{\Gamma}_{Z^{\prime}}(\mathcal{O}_{\mathcal{P}^{\prime}, \mathbb{Q}}({}^{\dagger}H^{\prime}))[c]\),
initially a complex of
\(\mathscr{D}^{\dagger}_{\mathcal{P}^{\prime},\mathbb{Q}}({}^{\dagger}H^{\prime})\)-modules,
is concentrated in degree \(0\) only.  Consequently, by
\cite[Theorem~2.3]{baldassarri-berthelot:dwork-cohomology-for-singular-hypersurfaces}
(which asserts that the category of
\(\mathscr{D}^{\dagger}_{\mathcal{P}^{\prime},\mathbb{Q}}({}^{\dagger}H)\)-modules
is equivalent to the category of modules over the overconvergent Weyl
algebra, the equivalence being induced by the functor
\(\mathrm{H}^{0}(\mathcal{P}^{\prime},-)\)), we find that
\[
\mathrm{H}^{0}(\mathcal{P}^{\prime}; \mathbb{R}\underline{\Gamma}_{Z^{\prime}}(\mathcal{O}_{\mathcal{P}^{\prime}, \mathbb{Q}}({}^{\dagger}H^{\prime}))[c])
\]
is a module over the overconvergent Weyl algebra
\(\mathrm{H}^{0}(\mathcal{P}^{\prime};\mathscr{D}^{\dagger}_{\mathcal{P}^{\prime},\mathbb{Q}}({}^{\dagger}H^{\prime}))\),
and is thus concentrated in degree \(0\) only.

\medskip%
Now, the zeroth row of \eqref{eq:smaller}, shifted to the left by \(r\), is
\[
{B} \to B^{\oplus c} \to B^{\oplus \binom{c}{2}} \to \cdots \to B^{\oplus \binom{c}{c-1}} \to \underset{\bullet}{B^{\oplus \binom{c}{c}}}
\]
(where the bullet indicates the degree zero item of the chain
complex).  We want to show it is acyclic except in degree \(0\).  Since this is
a relative de~Rham complex,  according to
\cite[p.~197, Remark]{baldassarri-berthelot:dwork-cohomology-for-singular-hypersurfaces},
it represents the complex
\(\mathrm{H}^{0}(\mathcal{P}^{\prime}; \varpi^{\prime}_{+}\mathcal{L})\),
which is isomorphic to
\[
\mathrm{H}^{0}(\mathcal{P}^{\prime}; \mathbb{R} \underline{\Gamma}_{Z^{\prime}} (\mathcal{O}_{\mathcal{P}^{\prime}, \mathbb{Q}} ({}^{\dagger} H^{\prime})) [c])
\]
by Theorem~\ref{theorem:local-comparison}.  As discussed in the
preceding paragraph, this latter complex is indeed concentrated in
degree \(0\).  This completes the proof.
\end{proof}

\begin{thebibliography}{99}
\bibitem[8]{baldassarri-berthelot:dwork-cohomology-for-singular-hypersurfaces} Francesco Baldassarri and Pierre Berthelot. On {D}work cohomology for singular hypersurfaces. In {\em Geometric aspects of {D}work theory. {V}ol. {I}, {II}}, pages 177--244. Walter de Gruyter, Berlin, 2004.
\bibitem[11]{berthelot:arithmetic-d-modules-1-differential-operators-of-finite-level} Pierre Berthelot. {D}-modules arith\'etiques. {I}. {O}p\'erateurs diff\'erentiels de niveau fini. {\em Ann. Sci. \'Ecole Norm. Sup. (4)}, 29(2):185--272, 1996.
\end{thebibliography}
\end{document}
